\section{Learning Theories}\label{sec:id-theories}
\owner{Group}
\pages{2}

Four learning theories shape the app. Each one is applied to a specific design decision, so that the link between theory and design can be checked.

\subsection{Cognitive Load Theory}
Cognitive load theory holds that working memory has a small capacity, and that instruction fails when a learner must hold too many new ideas at once \cite{Sweller1988}. A BST lesson can overload a novice quickly, because the ordering rule, the comparison, the choice of subtree and the position of the node all need to be tracked together. The app applies the theory in three ways. Each trace shows one comparison at a time, the step controls let the student pause on any step, and the explanation for each step names one rule rather than several. The worked-example effect \cite{Sweller1988} is applied in the first lesson, where a fully worked search is shown before the student attempts one.

\subsection{Scaffolding and the Zone of Proximal Development}
Scaffolding is support that is given while a learner works on a task they cannot yet complete alone, and that is withdrawn as competence grows \cite{Wood1976}. The idea is related to Vygotsky's zone of proximal development \cite{Vygotsky1978}, which describes the range of tasks a learner can do with help but not yet alone. The app applies scaffolding in two places. The diagnostic gate routes a student to the lesson that matches a specific misconception, rather than to all the material. Inside a trace, the three hint levels give progressively more help, and a student who answers correctly without a hint moves on with less support.

\subsection{Retrieval Practice}
Retrieving information from memory strengthens later recall more than rereading the same material does \cite{Roediger2006}. The app uses retrieval in two ways. The diagnostic asks the student to classify and predict before any explanation is given. During a trace the student predicts the next step before the app shows it. In both cases the student must produce an answer, and the feedback follows the attempt rather than preceding it.

\subsection{Mastery Learning}
Mastery learning holds that most students can reach a defined standard when they are given the time and feedback they need \cite{Bloom1968}. The app applies this by treating each objective as a standard to be reached, not as a single score to be passed once. A student who has not reached the standard for an objective is routed back to the lesson for it, and the progress record shows the objective as practised or secure, as set out in Section~\ref{sec:id-progress}.
